\chapter{Pensar como Quem Testa um Pedido}
\label{ch:mentalidade}
Um colega pede-te o prompt que usaste ontem. Tens o texto, mas já não tens a mensagem original nem sabes qual das quatro respostas foi enviada. Conseguirias repetir o teste? Guardar só a frase que escrevemos é como anotar a receita e deitar fora a lista de ingredientes.

\section{Tratar o pedido como uma especificação}
Um critério de aceitação diz o que tem de acontecer para podermos usar a resposta. Na ficha de R-014, uma resposta aceitável informa que o pedido foi recebido e indica que o levantamento fica por confirmar. Se a ficha passar a «pronto para levantamento», o aviso permitido muda. Tens dois casos que exercitam lados diferentes da mesma regra.

A comparação com código ajuda na disciplina. Guarda versões e casos de teste, anota erros que consegues reproduzir e regista por que mudaste o pedido. Ao executar código, a linguagem define operações precisas. Num modelo de linguagem, a resposta depende do modelo, do contexto e da configuração. A especificação diz o que aceitamos da saída. Os testes mostram como se comportou nos exemplos avaliados.

Podes fazer este percurso num chat. Num programa que trata vários pedidos, guarda a regra reutilizável separada das fichas de cada caso. O capítulo~\ref{ch:personas} mostra como a aplicação pode fazer essa separação.

\begin{promptlabbox}{Reescrever um pedido solto como “função”}
\textbf{Objetivo e contexto:} Tornar verificável um pedido de resumo. Usa um texto fictício ou autorizado e define os factos essenciais.

Pega num pedido típico que fazes a um LLM (por exemplo, “faz um resumo deste texto para o meu chefe”) e reescreve-o em formato quase de função:

\medskip

\begin{lstlisting}
Função: RESUMO_EXECUTIVO
Papel do modelo: Assistente de comunicação clara, concisa e pragmática.
Entrada:

Texto original (até 2000 palavras)

Perfil do destinatário (ex: "gestor ocupado", "equipa técnica")
Saída:

Resumo em 5 a 7 frases curtas

Tom: profissional, direto, sem jargão desnecessário

Destaques: riscos, oportunidades, decisões a tomar

Instruções:

Não inventar factos que não estejam no texto original.

Se o texto for demasiado vago, pedir esclarecimentos em vez de adivinhar.
\end{lstlisting}

Experimenta o mesmo texto e destinatário com o pedido original e com a versão organizada.

\textbf{Resultado esperado e validação:} Verifica comprimento, público, cobertura dos factos e informação acrescentada. Assinala o trecho do texto de origem que sustenta cada decisão mencionada.

\textbf{Falhas possíveis:} Não fornecer o texto, alterar o público entre condições ou presumir que um formato de função cria execução determinística.
\end{promptlabbox}

\section{Quatro maneiras de perder o controlo da tarefa}

\subsection*{1. Pedir uma recomendação sem critérios}

Escrever perguntas como se o modelo fosse um motor de pesquisa:

\begin{quote}
\small
«Qual é a melhor ferramenta para construir este site?»
\end{quote}

A pergunta omite o tipo de aplicação, a equipa e os critérios de escolha. Antes de procurar nomes, decide que informação te permitiria compará-los.

\subsection*{2. Empilhar tarefas}

Misturar tarefas diferentes numa sopa única:

\begin{quote}
\small
«Resume estas mensagens, responde a todas, prepara a lista de decisões e escreve o relatório da semana.»
\end{quote}

Se o resumo omitir um pedido importante, a resposta e o relatório podem repetir a omissão. Separa o trabalho e confirma primeiro a lista de pedidos.

\subsection*{3. Falta de exemplos quando o formato importa}

Pedir algo complexo sem mostrar um exemplo:

\begin{quote}
\small
“Organiza estes e-mails numa tabela com colunas úteis.”
\end{quote}

Indica quem vai usar a tabela, para que decisão e o que deve aparecer em cada coluna. Uma linha preenchida mostra a diferença entre «data do pedido» e «data em que podemos responder».

\subsection*{4. Confiança cega}

Se a resposta contém um valor, localiza o registo que o sustenta. Se apresenta uma referência, abre a fonte e verifica se ela diz aquilo que a frase afirma. Outra resposta parecida do mesmo modelo pode repetir o erro inicial.

\section{Ciclo básico: protótipo \texorpdfstring{$\to$}{->} teste \texorpdfstring{$\to$}{->} refinamento \texorpdfstring{$\to$}{->} versão estável}

Para qualquer tarefa recorrente, pensa num mini-ciclo de desenvolvimento de prompt:

\subsection*{1. Protótipo rápido}

Escreve a tua melhor tentativa inicial. Não te preocupes com perfeição; preocupa-te com clareza mínima:

\begin{itemize}[leftmargin=*]
\item Define a tarefa e, se ajudar, o papel do modelo.
\item Especifica a entrada e a saída.
\item Indica restrições básicas (comprimento, tom, formato).
\end{itemize}

\subsection*{2. Teste com casos representativos}

Aplica o pedido a casos fictícios realistas ou a dados cuja utilização esteja autorizada. Inclui:

\begin{itemize}[leftmargin=*]
\item casos fáceis,
\item casos ambíguos,
\item casos extremos em que o sistema pode falhar.
\end{itemize}

\subsection*{3. Ajuste com base em erros concretos}

Quando o modelo falhar, não culpes “a IA” em abstrato. Pergunta:

\begin{itemize}[leftmargin=*]
\item Que instrução faltou?
\item Que exemplo podia ter evitado este erro?
\item Ficou claro o que a resposta deve entregar?
\end{itemize}

Depois, ajusta o prompt como ajustarias código após um bug report.

\subsection*{4. Normalização}

Quando a versão cumprir os critérios dos testes definidos, guarda uma referência para voltar a ela.

\begin{itemize}[leftmargin=*]
\item Dá-lhe um nome (“Resumo-Executivo-v1.2”).
\item Guarda-o num sítio acessível (repositório, wiki, código).
\item Documenta o que faz, o que não faz e exemplos de uso.
\end{itemize}

Mais à frente, LLMOps designa a gestão e operação de sistemas com modelos (capítulo~\ref{ch:llmops}); os guardrails são mecanismos de validação e controlo. Por agora, guarda o pedido, os resultados e a razão de cada mudança. É esse registo que permite aprender com as falhas.

\begin{promptlabbox}{Aplicar o ciclo a uma tarefa tua}
\textbf{Objetivo e contexto:} Desenvolver uma versão de prompt com entradas representativas e outras reservadas para teste.

Escolhe uma tarefa que fazes \emph{muitas} vezes com LLMs (por exemplo, traduzir e-mails técnicos, gerar descrições de produtos, resumir reuniões).

\medskip

Passos:

\begin{enumerate}[leftmargin=*]
\item Escreve uma primeira versão do pedido, com a entrada e a saída esperada.
\item Testa com cinco casos fictícios representativos ou autorizados.
\item Para cada falha, guarda a entrada e a razão da alteração. Avalia a nova versão em casos reservados antes de começar.
\item Quando estiveres satisfeito, dá um nome de versão e guarda.
\end{enumerate}

Agora tens uma versão e evidência do que funcionou. Ainda falta testar casos que não usaste no refinamento.

\textbf{Resultado esperado e validação:} Guarda versão, falha, alteração e resultado. Reavalia os casos anteriores e os reservados; não declares generalização a partir dos exemplos usados na revisão.

\textbf{Falhas possíveis:} Ajustar o prompt até decorar o teste ou apagar resultados negativos do registo.
\end{promptlabbox}

\section{Antes de voltar a usar um pedido}

\begin{checklistbox}{Checklist rápida de qualidade de prompt}
Antes de usar um prompt de forma recorrente, verifica:

\begin{itemize}[leftmargin=*]
\item \textbf{Papel, quando útil}: a perspetiva pedida ajuda a tarefa sem fingir credenciais?
\item \textbf{Objetivo definido}: o que vai a resposta permitir decidir ou fazer?
\item \textbf{Formato de saída}: está especificado (lista, JSON, tabela conceptual, texto corrido)?
\item \textbf{Exemplos}: há uma entrada e uma resposta esperada para um caso difícil?
\item \textbf{Limites}: está dito o que o modelo \emph{não} deve fazer (inventar dados, dar opiniões legais, etc.)?
\item \textbf{Gestão de incerteza}: está explicado como agir se a informação for insuficiente (pedir mais dados, indicar “não sei”)?
\item \textbf{Linguagem}: está definido quem vai usar a resposta e o que já conhece?
\end{itemize}
\end{checklistbox}

Escolhe os critérios que alteram a aceitação da resposta. Para um título, o tom pode ser decisivo. Para extrair uma data, interessa poder encontrá-la no documento.

\section{O caso que fica para o teste}
Reserva alguns exemplos antes de começares a afinar o pedido. Se uma falha num caso reservado te levar a mudar o prompt, esse caso passa a fazer parte do desenvolvimento. Usa casos novos para voltar a avaliar. Cinco mensagens da mesma pessoa com a mesma dúvida dizem pouco sobre situações diferentes.

Na biblioteca, uma mensagem pode pedir ao mesmo tempo o estado de uma reserva e a renovação de outro empréstimo. No próximo capítulo, esse caso obriga-nos a explicar a regra com um exemplo.
